% Generated by tex/build_swebok_structure.py from tex/chapters/ch01/03_ソフトウェア要求の基礎.tex
\subsubsection{機能要求}

\term{機能要求}とは、ソフトウェアが提供すべき観測可能な振る舞いを表す要求である。言い換えると、「ソフトウェアが何をするか」を述べる要求である。

機能要求は、二つの観点で理解しやすい。第一に、\term{方針}である。方針とは、常に守られるべき規則である。たとえば銀行システムでは、「口座は少なくとも一人の所有者を持つ」「口座残高は負にならない」といった要求が該当する。第二に、\term{プロセス}である。プロセスとは、実行される業務手続きである。入金、出金、振込、予約、取消、検索、承認などが該当する。

\begin{definitionbox}機能要求とは、ソフトウェアが実行すべき処理、守るべき業務規則、利用者に提供すべき振る舞いを述べる要求である。\end{definitionbox}
